\section{Introduction}

%FELA in general
% 2.0: FELA introduction and we only focus on LB in this paper.

%The determination of the ultimate load-carrying capacity, or collapse load, of geotechnical structures is a fundamental task in stability analysis and design. Finite Element Limit Analysis (FELA) provides a rigorous numerical framework for this purpose by computing strict upper and lower bounds to the true collapse load~\cite{Sloan_stability_analysis, Dunne_thesis}. This paper focuses on lower-bound FELA, where the objective is to construct a statically admissible stress field that satisfies equilibrium, boundary conditions, and yield admissibility.


% Convex failure criteria, converts the problem to cvxopt problems.
% 2.0: LB FELA in 2D and 3D leads to large-scale SOCPs and even SDPs.

%For many convex yield criteria, the lower-bound FELA formulation can be cast as a conic optimization problem. In 2D, criteria such as Tresca and Mohr-Coulomb can be represented using second-order cone (SOC) constraints, leading to Second-Order Cone Programming (SOCP) formulations~\cite{2DLB, UB_simplex}. Similar SOCP structures also arise in selected three-dimensional formulations, for example with von Mises or Drucker-Prager criteria, while other yield criteria may require semidefinite or nonlinear programming representations~\cite{3DLB, 3DLD, 3DHB, 3DBishop}. When Bernstein polynomial stress elements and adaptive mesh refinement are employed~\cite{Bernstein_FELA}, the resulting discretized problems involve solving optimization problems with an increasing number of linear equilibrium constraints and local yield constraints. For the Drucker-Prager criterion, these yield constraints can be represented as SOC constraints, producing large-scale SOCPs that demand scalable numerical solvers.

% IPM used to be the solver of choice for the past 2 decades.

%For the past decades, the Interior Point Method (IPM) has been the method of choice to solve these problems. However, the IPM theoretically suffers from scalability issues because it needs to reformulate, and refactorize the Karush-Kuhn-Tucker (KKT) matrix in every iteration. Although the development of bespoke, high-performance IPM solvers—exploiting problem structure and GPU parallelism—has alleviated some of these challenges, and such solvers may outperform the state-of-the-art commercial solver MOSEK~\cite{MOSEK, MOSEK_PD, high_performance_IPM} in certain cases, the fundamental limitations of IPMs remain significant in large-scale applications. 

%In the 2010s, several efficient conic solvers were developed for SOCPs based on primal-dual interior-point methods. CVXOPT adopts a dense linear algebra approach using self-dual embedding and Cholesky factorization of the reduced KKT system~\cite{CVXOPT}, while ECOS employs Mehrotra’s predictor-corrector strategy with Nesterov–Todd scaling, using sparse $LDL^\top$ factorization and fixed symbolic ordering for better performance on embedded systems and small problems~\cite{ECOS}.

% Primal-dual IPM algorithms.

%Among some most recent implementations of primal-dual interior point methods, Clarabel is an open-source general purpose solver that solves conic optimization problems, with supported cone types including the second-order cone and exponential cone. It incorporated recent algorithmic advancements including homogeneous embedding and Nesterov-Todd (NT) scaling, allowing it to outperform existing IPM solvers in a series of benchmarks~\cite{Clarabel}. It also comes with a GPU implementation, CuClarabel, utilizing the SIMD (single instruction, multiple data) structure of the GPU to accelerate their IPM method~\cite{CuClarabel}. 


%In recent years, tailored solvers designed for FELA have also been developed. Podlich developed a tailored high-performance interior-point solver that explicitly forms and factorizes the Schur complement system arising in IPM~\cite{high_performance_IPM}. To reduce fill-in and improve numerical efficiency, effective presolve techniques such  as matrix reordering techniques and variable elimination, a modified Schur complement method for handling dense columns, and a specialized routine for eliminating fixed variables under cone constraints~\cite{Podlich_2017}. These structural enhancements, resulted in a 36\% reduction in runtime and fewer nonzeros in the factorized matrices. Combined with parallelization and GPU acceleration, the solver achieved up to a 4.5× speedup over the commercial solver MOSEK in 3D FELA problems with Drucker-Prager criterion, resulting a problem in SOCP form.


% 2.0: IPM-based solvers: mature, strong, but bottlenecked by factorization/memory

Interior-point methods (IPMs) have long been the standard approach for solving such conic optimization problems, owing to their robustness, high accuracy, and fast convergence in terms of iteration count. Over the past two decades, both general-purpose conic solvers and tailored FELA solvers have been developed around this methodology. General-purpose solvers such as CVXOPT, ECOS, MOSEK, and Clarabel provide mature, robust primal-dual interior-point implementations for SOCPs, while tailored FELA solvers further exploit problem-specific structures, including Schur-complement systems, variable elimination, matrix reordering, and parallel factorization~\cite{CVXOPT, ECOS, MOSEK, Clarabel, high_performance_IPM, Podlich_2017}. These solvers vary in their presolve techniques, problem formulations, and linear-solve strategies, yet they commonly rely on repeatedly forming and factorizing iteration-dependent KKT systems. As the problem size increases, particularly in 3D finite element analyses, the resulting factorization cost and fill-in can lead to substantial memory demands, making sparse direct methods difficult to apply on very large instances~\cite{IPM25}. Indirect or iterative linear solvers are therefore an attractive alternative for reducing memory consumption~\cite{Saad_book}. However, the linear systems arising in IPMs often become increasingly ill-conditioned as the iterates approach convergence, which can significantly degrade the effectiveness of indirect methods~\cite{GeNIOS}, though that problem may be alleviated by the use of carefully designed preconditioners~\cite{Matrix_free_IPM, Gondzio_2023}.


% 2.0: ADMM opportunity and limitation




Recent advances have seen the application of a first-order optimization method, ADMM to FELA~\cite{daSilva_2020, Antao_2021}, demonstrating its potential as a robust solver that can yield answers with good accuracy in a reasonable amount of time. Compared with IPMs, ADMM offers a different computational trade-off: it typically requires more iterations, but each iteration can be significantly cheaper. It uses only first-order derivative information, bypassing the need to reformulate the Hessian matrix every iteration. Therefore ADMM is able to store and reuse the results obtained from the initial numerical factorization, thereby benefiting from cheaper per-iteration computational cost and low memory requirement. 



Another attractive feature of ADMM in FELA is its projection-centric treatment of yield constraints. Previous studies have reported successful applications of ADMM to more sophisticated 3D yield criteria, such as Lade-Duncan and Hoek-Brown criteria~\cite{3DLD, 3DHB}. This projection-centric treatment provides a practical alternative to interior-point formulations that require the yield locus to admit an explicit conic (SOCP/SDP) representation, which can be restrictive for more sophisticated criteria. Although the present work focuses on SOCP formulations, this flexibility is important for future extensions to more general yield models that are more common in geotechnical practice.
 

Despite these advantages, ADMM typically requires substantially more iterations than IPMs to reach a high-accuracy solution~\cite{daSilva_2020, ADMM_Book}, and its convergence behaviour is strongly influenced by the choice of penalty parameters in the augmented Lagrangian. General-purpose operator-splitting solvers address this issue in different ways. For example, OSQP employs adaptive penalty updates to balance primal and dual residuals~\cite{OSQP}, while SuperADMM introduces constraint-wise penalty adaptation to accelerate convergence~\cite{SuperADMM}. SCS supports both direct and matrix-free indirect linear-system solves. Although general-purpose ADMM solvers such as OSQP, COSMO, and SCS have demonstrated the effectiveness of operator-splitting methods for large-scale quadratic programming (QP) and conic optimization problems~\cite{OSQP, COSMO, SCS}, they are not specifically designed to exploit the structure of lower-bound FELA. 

Lower-bound FELA possesses several features that make a more specialized treatment attractive. Its optimization problems combine global sparse equilibrium constraints with a large number of small, local yield constraints that share a common conic structure. At the same time, adaptive mesh refinement naturally generates a sequence of closely related optimization problems, while the linear systems arising within ADMM are repeatedly solved with closely related or unchanged sparsity patterns. These structures appear at different levels of the computation and provide distinct opportunities for acceleration.



The central idea of this work is therefore to exploit these complementary structures jointly within a single large-scale solution framework. Repeated ADMM linear systems permit expensive sparse factorizations to be reused across iterations; the local structure of the yield constraints enables massively parallel cone projections and selective active-cone penalty reweighting; and successive mesh refinements provide a natural hierarchy for coarse-to-fine warm starting. These algorithmic features are combined with GPU-based sparse linear algebra and customized CUDA kernels to reduce the cost of the dominant operations. Building on this structure-aware formulation, we develop a GPU-accelerated ADMM framework for large-scale SOCPs arising in two- and three-dimensional lower-bound FELA.

The main contributions of this work are:

\begin{enumerate}
  \item We derive a tailored ADMM formulation for the resulting SOCPs, including the condensed linear system, cone projection steps, termination criteria, and two optional enhancements: ellipsoidal projection and an augmented linear-system formulation.
  \item We propose two acceleration strategies for large-scale adaptive FELA: a coarse-to-fine warm-start strategy and an active-cone based dynamic penalty reweighting scheme inspired by SuperADMM~\cite{SuperADMM}.
  \item We implement the dominant computational kernels on the GPU, including sparse direct solves using cuDSS~\cite{CUDSS}, customized vector operations and cone projections, and cuBLAS/cuSPARSE-based matrix operations, thereby minimizing CPU-GPU data transfer.
  \item We evaluate the proposed framework on representative 2D and 3D lower-bound FELA benchmarks, including rigid platen, strip footing, and vertical cut problems, and discuss the distinct computational bottlenecks and speedup mechanisms in two and three dimensions.
\end{enumerate}

